\documentclass[aspectratio=169, 10pt]{beamer}

\usepackage{xeCJK}
\setCJKmainfont{Microsoft JhengHei}
\setCJKsansfont{Microsoft JhengHei}
\setCJKmonofont{KaiTi}

\usepackage{booktabs}
\usepackage{multicol}
\usepackage{amsmath, amssymb}
\usepackage{graphicx}
\usepackage{tikz}

\usetheme{ntu}
% 設定台大藍主色與輔助色
\setcolor[ntugold]{ntublue}

\setdepartment{國立臺灣大學 工學院 工程科學及海洋工程學系 iCAN 智慧計算與網路實驗室}

\title{應用知識圖注意力網路於食譜推薦之\\效能評估與可解釋性研究}
\subtitle{碩士學位論文口試}
\author{研究生：廖健合 \quad (指導教授：張瑞益 教授)}
\date{\today}

\begin{document}

\inserttitlepage

\begin{frame}{目錄 Outline}
  \tableofcontents
\end{frame}

\section{研究背景與動機}
\begin{frame}{食譜平台的資訊過載與推薦需求}
  \begin{mybox}{ntublue}{背景與痛點}
    \begin{itemize}
      \item 食譜平台累積海量食譜、食材與使用者互動。
      \item 純協同過濾 (CF) 受限於資料稀疏與冷啟動問題。
    \end{itemize}
  \end{mybox}
  
  \begin{examplebox}{核心研究目標}
    如何有效利用食材與標籤等輔助知識，產生高準確且可說明的推薦結果？
  \end{examplebox}
\end{frame}

\end{document}
